% =============================================================================
% BIS 527: Introduction to Health Data Science
% SCRIBE NOTES TEMPLATE
%
% -----------------------------------------------------------------------------
% NEVER USED LATEX BEFORE? START HERE.
% -----------------------------------------------------------------------------
% LaTeX is a typesetting system. You write a plain text file (this one, ending
% in .tex), then a program turns it into a PDF. You do not see the formatting
% while you type, the way you would in Word. You describe what things ARE
% ("this is a section heading") and LaTeX decides how they LOOK.
%
% Three things to know before you touch anything:
%
%   1. A line starting with % is a COMMENT. LaTeX ignores it. Everything you
%      are reading now is a comment. You can delete all of it.
%
%   2. A word starting with a backslash is a COMMAND, for example \section.
%      Text in curly braces {} after a command is what the command acts on.
%      So \section{Motivation} makes a section heading that reads "Motivation".
%
%   3. A blank line starts a new paragraph. One blank line is enough. Extra
%      blank lines change nothing. A line break WITHOUT a blank line does not
%      start a new paragraph, so you can wrap long lines however you like.
%
% Special characters. These have meaning in LaTeX, so to print them literally
% you put a backslash first:  \%  \$  \&  \#  \_  \{  \}
% Quotation marks are two backticks and two apostrophes: ``like this''
% A dash between numbers (a range) is two hyphens: 1908--1922
%
% -----------------------------------------------------------------------------
% HOW TO GET A PDF OUT OF THIS FILE
% -----------------------------------------------------------------------------
% EASIEST: use Overleaf (overleaf.com). Make a free account, create a project,
% upload this file and scribe-refs.bib, and press Recompile. Nothing to install.
%
% ON YOUR OWN MACHINE: install TeX Live (or MacTeX on a Mac), then run these
% four commands in the folder that holds this file. The order matters.
%
%     pdflatex scribe-lec03-yourname
%     bibtex   scribe-lec03-yourname
%     pdflatex scribe-lec03-yourname
%     pdflatex scribe-lec03-yourname
%
% Why four? The first pass writes down which sources you cited. bibtex looks
% them up in scribe-refs.bib. The last two passes put the citations and the
% cross-references in place. If your citations show up as [?], you skipped a
% pass. Run all four again.
%
% -----------------------------------------------------------------------------
% WHAT TO DO WITH THIS FILE
% -----------------------------------------------------------------------------
%   1. Copy it and rename the copy   scribe-lec<NN>-<yourlastname>.tex
%      for example  scribe-lec03-nguyen.tex
%   2. Fill in the four \renewcommand lines just after \begin{document}.
%   3. Write your chapter in the sections below.
%   4. Submit BOTH the .tex source and the compiled .pdf on Canvas.
%
% You will also need scribe-refs.bib, the shared bibliography file, in the same
% folder. It is on the course website.
%
% -----------------------------------------------------------------------------
% WHAT A SCRIBE NOTE IS
% -----------------------------------------------------------------------------
% Each note is an INDEPENDENT TEXTBOOK CHAPTER. A reader who attended no
% lecture and read no other chapter must be able to follow it start to finish.
% It is not a transcript.
%
%   - Define every symbol you use, in this chapter, even if another chapter
%     already defined it. Repetition across chapters is correct.
%   - Do not write "as we saw in Lecture 2" in place of content. State the
%     result again, in full.
%   - Say at the start what background you assume, if you assume any.
%
% LENGTH: 7 to 8 pages, roughly 2,400 words including tables and figures.
% This is a limit, not a target. A shorter chapter that covers the material
% beats a longer one. Three habits blow the budget: stating an idea, then
% restating it as an aside, then again as an example; a worked example that
% repeats arithmetic done earlier; and a long list where only two entries
% carry the argument.
%
% -----------------------------------------------------------------------------
% HOW TO WRITE IT
% -----------------------------------------------------------------------------
%   1. NO ANTHROPOMORPHIZATION. An algorithm does not know, want, decide, or
%      look at anything. Neither does a field, a method, or a model. Write
%      "the procedure compares A[mid] with t", not "the procedure looks at the
%      middle element". Write "a researcher draws on statistics for the limits
%      of what the observations support", not "statistics tells us the limits".
%      A thing may act mechanically (a step applies, a procedure halts, a rule
%      assigns treatment). An abstraction may not perform a linguistic, formal,
%      or evaluative act (a word does not name, a model does not formalize, a
%      lineage does not judge).
%
%   2. ONE CLAIM PER SENTENCE. Do not join several arguments with semicolons.
%
%   3. EVERY SENTENCE MUST ASSERT SOMETHING. Cut sentences that only announce
%      what is coming. "Take the compound at face value." "Two short histories
%      lie behind data science." "The commitment was not always kept." Each of
%      those makes the reader wait for the next sentence to learn anything.
%      Open each paragraph on its claim.
%
%   4. NO RHETORIC. Cut "genuinely", "remarkably", "the whole game", "notice
%      that", "habits of mind", and second-person flourishes.
%
%   5. NO NOTATION OR JARGON WITHOUT A DEFINITION FIRST. If you write $\Theta$,
%      define $\Theta$. If you write "stable sort", define it.
%
%   6. MOTIVATION BEFORE FORMALISM, and a concrete example before the general
%      statement. A reader new to the material must be able to follow it.
%
% Use bolded run-in headers (\paragraph{Setup.}) to signal structure inside a
% section, and \emph{} for emphasis. Declare symbols in prose, and use := for
% definitions.
%
% -----------------------------------------------------------------------------
% THE SECTION SKELETON
% -----------------------------------------------------------------------------
% The skeleton below is the default, so that the notes read as one book rather
% than twelve separate documents. It is a suggestion, not a rule. Depart from
% it where the material does not fit it, and say why in a comment at the top of
% your file. The instructor example on the course website departs from it in
% four places and explains each one.
%
% -----------------------------------------------------------------------------
% NOTHING IN THIS COURSE IS HANDWRITTEN. That includes figures: use TikZ, or
% plot in Python and include the image file. Photographs of written work are
% not accepted. Board diagrams from lecture must be redrawn, not photographed.
% =============================================================================


% =============================================================================
% PREAMBLE
% -----------------------------------------------------------------------------
% Everything between here and \begin{document} sets up the formatting. You do
% not need to change any of it. Each line is explained so you know what it
% does, not because you have to edit it.
% =============================================================================

% The document class. "article" is the standard class for short documents.
% 12pt is the font size and letterpaper is the page size, both required.
\documentclass[12pt,letterpaper]{article}

% A "package" adds features. \usepackage loads one.

\usepackage[margin=1in]{geometry}   % page margins, 1 inch on all four sides
\usepackage{setspace}               % control over line spacing
\onehalfspacing                     % one-and-a-half spacing, required

% Fonts and mathematics. The order of these five lines matters: newtxmath
% redefines two commands that amssymb and amsthm also define, so loading them
% in a different order gives an error. Leave this block alone.
\usepackage[T1]{fontenc}            % modern font encoding, fixes accented letters
\usepackage{amsmath}                % the standard mathematics package
\usepackage{amssymb}                % extra mathematical symbols
\let\Bbbk\relax                     % clears a name clash before newtxmath
\usepackage{newtxtext}              % Times text font
\usepackage{newtxmath}              % Times mathematics font
\let\openbox\relax                  % clears a second name clash
\usepackage{amsthm}                 % theorem, definition and proof environments

\usepackage{microtype}              % small spacing improvements, purely cosmetic
\usepackage{graphicx}               % lets you include image files
\usepackage{booktabs}               % better horizontal rules in tables
\usepackage{xcolor}                 % colours, used by the links and code style
\usepackage{listings}               % typesetting source code
\usepackage{algorithm}              % a float for pseudocode
\usepackage{algpseudocode}          % the pseudocode commands themselves
\usepackage{caption}                % control over caption formatting
\usepackage{natbib}                 % author-year citations, \citet and \citep
\usepackage{etoolbox}               % used by the three \AtBeginEnvironment lines

% The 1.5 spacing requirement is for body text. Reference lists, tables and
% code are read as blocks, not as running prose, so they are set tight.
\AtBeginEnvironment{thebibliography}{\small\setstretch{1.0}}
\setlength{\bibsep}{2pt}
\AtBeginEnvironment{algorithmic}{\setstretch{1.0}}
\AtBeginEnvironment{tabular}{\setstretch{1.0}}
\captionsetup{font={small,stretch=1.0}, skip=4pt}

\usepackage{hyperref}               % makes citations and references clickable
\usepackage{titlesec}               % control over heading size and spacing

\hypersetup{colorlinks=true, linkcolor=blue!45!black,
            citecolor=blue!45!black, urlcolor=blue!45!black,
            pdflang={en-US}}

% Heading sizes and the space above and below them.
\titleformat{\section}{\large\bfseries}{\thesection}{0.7em}{}
\titleformat{\subsection}{\normalsize\bfseries}{\thesubsection}{0.7em}{}
\titlespacing*{\section}{0pt}{1.0\baselineskip}{0.3\baselineskip}
\titlespacing*{\subsection}{0pt}{0.8\baselineskip}{0.25\baselineskip}

% --- The title block ---------------------------------------------------------
% \newcommand defines a new command. These four are placeholders. You do not
% edit them here; you override them after \begin{document}, which is where a
% reader of your file will look for them.
\newcommand{\lecturenumber}{N}
\newcommand{\lecturetitle}{Title of This Lecture}
\newcommand{\scribename}{Your Name}
\newcommand{\scribenetid}{abc12}

% This command draws the title block using the four values above. The % at the
% end of some lines is not a comment marker here; it stops LaTeX from inserting
% an unwanted space. Leave them in.
\newcommand{\makescribetitle}{%
  \begin{center}
    {\LARGE\bfseries Lecture \lecturenumber: \lecturetitle\par}
    \vspace{1.2em}
    {\large \scribename\par}
    \vspace{0.2em}
    {\normalsize \texttt{\scribenetid @yale.edu}\par}
    \vspace{0.9em}
    {\normalsize BIS 527: Introduction to Health Data Science\\
     Yale School of Public Health\\
     Lecturer: Harsh Parikh \quad\textbullet\quad Fall 2026\par}
  \end{center}
  \vspace{0.8em}
}

% --- The TL;DR block ---------------------------------------------------------
% \newenvironment defines a new environment, that is, a pair of commands you
% write as \begin{name} ... \end{name}. This one sets your summary in a
% narrower measure so it reads as a summary and not as the first paragraph.
\newenvironment{tldr}{%
  \begin{center}
    \bfseries\large TL;DR
  \end{center}
  \vspace{-0.4em}
  \begin{quote}\setstretch{1.25}\normalsize\noindent\ignorespaces
}{%
  \end{quote}\vspace{0.8em}
}

% --- The AI use declaration --------------------------------------------------
% Draws a box around whatever you pass it. Course policy is below, at the point
% where you fill it in.
\newcommand{\aideclaration}[1]{%
  \begin{center}
  \fbox{\parbox{0.93\textwidth}{\setstretch{1.15}\small
    \textbf{AI use declaration.} #1}}
  \end{center}
  \vspace*{0.8em}
}

% --- Numbered environments for formal statements -----------------------------
% These give you \begin{definition} ... \end{definition} and so on. LaTeX
% numbers them for you and keeps the numbering consistent if you move things.
%
% Definitions, examples and remarks are set in upright type. Theorems and the
% things that behave like theorems are set in italic, which is the usual
% convention. Lemma, corollary and proposition share one counter with theorem,
% so you get Theorem 1, Lemma 2, Theorem 3 and never two things numbered 1.
%
% You do not need to declare a "proof" environment. amsthm supplies it, and it
% ends with a small hollow square, the QED symbol. See the example below.
\theoremstyle{definition}
\newtheorem{definition}{Definition}
\newtheorem{example}{Example}
\newtheorem{remark}{Remark}
\theoremstyle{plain}
\newtheorem{theorem}{Theorem}
\newtheorem{lemma}[theorem]{Lemma}
\newtheorem{corollary}[theorem]{Corollary}
\newtheorem{proposition}[theorem]{Proposition}

% --- How Python code is displayed --------------------------------------------
\lstdefinestyle{pythonstyle}{
  language=Python,
  basicstyle=\ttfamily\small\setstretch{1.0},
  keywordstyle=\color{blue!45!black}\bfseries,
  commentstyle=\color{gray}\itshape,
  stringstyle=\color{red!45!black},
  numbers=left, numberstyle=\tiny\color{gray},
  showstringspaces=false, breaklines=true,
  frame=single, framesep=3pt, tabsize=4
}
\lstset{style=pythonstyle}


% =============================================================================
% YOUR CHAPTER STARTS HERE
% =============================================================================
\begin{document}

% --- FILL IN THESE FOUR LINES ------------------------------------------------
% \renewcommand replaces the placeholder values set in the preamble.
% Keep the curly braces. Do not put a full stop inside them.
\renewcommand{\lecturenumber}{N}                  % e.g. 3
\renewcommand{\lecturetitle}{Title of This Lecture}% e.g. Sorting and Complexity
\renewcommand{\scribename}{Your Name}             % e.g. Alex Nguyen
\renewcommand{\scribenetid}{abc12}                % your NetID, no @yale.edu

\makescribetitle   % draws the title block from the four values above


% --- TL;DR -------------------------------------------------------------------
% 50 to 100 words. Say what the lecture established: the question it opened
% with, the main results or definitions, and what they are good for. Someone
% deciding whether to read your chapter should be able to tell from this block
% alone whether it covers what they need.
%
% Write this LAST. It is a summary of what you actually wrote, not a plan for
% what you intend to write.
\begin{tldr}
Replace this text with your summary.
\end{tldr}


% --- AI USE DECLARATION (REQUIRED) -------------------------------------------
% Course policy: you may use AI for anything, but you must declare the use, and
% auditing the output is your responsibility. An undeclared use, or an
% inaccuracy traced to one, scores zero.
%
% A real declaration has four parts:
%   1. Name the TOOL.
%   2. Say WHAT IT DID.
%   3. Say WHAT YOU VERIFIED, and how.
%   4. Say WHAT YOU DID NOT VERIFY.
%
% Part 4 is the one people leave out and the one that protects you. Claiming
% more checking than you did is worse than admitting the gap.
%
% If you used no AI tools, write:
%   No AI tools were used in preparing this note.
\aideclaration{Replace this text. Name the tool, say what it did, say what you
verified and how, and say what you did not verify.}


% =============================================================================
\section{Motivation}
% A concrete problem, before any formalism. State the question the lecture
% answers, in terms a reader outside the field would follow. Where the material
% is health-related, use a health example.


% =============================================================================
\section{Setup and Notation}

\paragraph{Notation.}
% Declare every symbol in prose before it is used. For example: "Let $A$ be an
% array of $n$ records, indexed $A[1], \dots, A[n]$. We write $t$ for the
% target value." Use := for definitions.


% =============================================================================
\section{Main Development}
% One idea per section. Give the informal statement, then the formal one, then
% the argument. Split this into as many sections as the material needs.


% =============================================================================
\section{Worked Example}
% Explicit numbers, carried through to a result. A reader should be able to
% reproduce every step. If the lecture used no numbers, delete this section and
% say so in a comment at the top of your file.


% =============================================================================
\section{Summary}
% What the chapter established, as a short list.


% =============================================================================
% SYNTAX EXAMPLES -- DELETE THIS ENTIRE SECTION BEFORE SUBMITTING
% -----------------------------------------------------------------------------
% Everything below shows you the syntax for the things you are most likely to
% need. Copy what you need into the sections above, then delete from this
% comment down to \end{document} and put the bibliography lines back.
% =============================================================================
\section*{Syntax examples (delete before submitting)}

\subsection*{Emphasis and quotation}
Use \emph{emphasis} for a term you are introducing. Quotation marks are
``two backticks and two apostrophes''. A number range uses two hyphens,
as in 1908--1922.

\subsection*{Citing a source}
Both commands pull from \texttt{scribe-refs.bib}. Use \verb|\citep| when the
citation sits in brackets at the end of a claim, and \verb|\citet| when the
author's name is part of your sentence.

Randomization separates an intervention from ordinary variation
\citep{fisher1935design}. \citet{tukey1962future} argued that data analysis
extends beyond mathematical statistics.

To cite a source that is not yet in \texttt{scribe-refs.bib}, open that file
and add an entry. Copy the shape of an entry that is already there. The key,
the first thing inside the braces, is what you type inside \verb|\citep{...}|.

\subsection*{Mathematics}
Mathematics inside a sentence goes between dollar signs, like $n \log n$.
Mathematics on its own line goes in an \texttt{equation} environment, which
numbers it so you can refer back to it:
\begin{equation}\label{eq:mean}
  \bar{x} := \frac{1}{n} \sum_{i=1}^{n} x_i .
\end{equation}
Refer to it with \verb|\eqref|, as in Equation~\eqref{eq:mean}. The tilde in
\verb|Equation~\eqref| is a space that will never break across a line.

\subsection*{A definition, and referring back to it}
\begin{definition}[Sorted array]
\label{def:sorted}
An array $A$ of length $n$ is \emph{sorted} if $A[i] \le A[i+1]$ for every
$i$ from $1$ to $n-1$.
\end{definition}

Definition~\ref{def:sorted} is what binary search requires. The pattern is
always the same: put a \verb|\label{...}| inside the thing you want to name,
then write \verb|\ref{...}| with the same text wherever you refer to it.
LaTeX fills in the number, so it stays correct when you move things around.
The same works for sections, tables, figures and equations.

\subsection*{A theorem and its proof}
State the claim in a \texttt{theorem} environment and argue it in a
\texttt{proof} environment. The proof ends with a hollow square, which LaTeX
adds for you. Everything inside a theorem is set in italic by convention, so
do not add italics yourself.

\begin{theorem}[Cost of binary search]
\label{thm:binary}
Let $A$ be a sorted array of $n \ge 1$ entries in the sense of
Definition~\ref{def:sorted}, and let $t$ be the target value. Binary search on
$A$ executes the body of its loop at most $\lfloor \log_2 n \rfloor + 1$ times.
\end{theorem}

\begin{proof}
Write $m_k$ for the number of entries still under consideration after $k$
iterations, so that $m_0 = n$. One iteration compares $A[mid]$ with $t$ and
then either returns or discards $A[mid]$ together with everything on one side
of it. Each of the two remaining pieces holds at most $\lfloor m_k / 2 \rfloor$
entries, so
\begin{equation}\label{eq:halving}
  m_{k+1} \le \left\lfloor \frac{m_k}{2} \right\rfloor
  \qquad\text{and therefore}\qquad
  m_k \le \left\lfloor \frac{n}{2^k} \right\rfloor .
\end{equation}
The loop runs only while at least one entry remains, so it stops at the first
$k$ with $\lfloor n / 2^k \rfloor = 0$. That inequality holds exactly when
$2^k > n$, which first occurs at $k = \lfloor \log_2 n \rfloor + 1$. The body
therefore runs at most $\lfloor \log_2 n \rfloor + 1$ times.
\end{proof}

A proof that ends on a displayed equation puts the square in the wrong place.
Write \verb|\qedhere| at the end of the display to fix it. A result you are
quoting rather than proving needs no proof environment, but it does need a
citation. \texttt{lemma}, \texttt{corollary} and \texttt{proposition} work the
same way as \texttt{theorem} and share its numbering, so a lemma following
Theorem~\ref{thm:binary} would be Lemma~2.

\subsection*{Lists}
An unnumbered list:
\begin{itemize}
  \item first point
  \item second point
\end{itemize}
A numbered list uses \texttt{enumerate} in place of \texttt{itemize}.

\subsection*{A table}
\begin{table}[h]
\centering
\begin{tabular}{lrr}
\toprule
Method & Comparisons & Time \\
\midrule
Linear search & $n$        & $O(n)$ \\
Binary search & $\log_2 n$ & $O(\log n)$ \\
\bottomrule
\end{tabular}
\caption{Put the caption below a table. The \texttt{lrr} says there are three
columns, the first left-aligned and the other two right-aligned.}
\label{tab:search}
\end{table}

\subsection*{A figure}
Save your plot as a PDF or PNG in the same folder, then write this. The
\texttt{width} setting scales the image to seven tenths of the text width.

\begin{verbatim}
\begin{figure}[h]
  \centering
  \includegraphics[width=0.7\textwidth]{myplot.pdf}
  \caption{Put the caption below a figure. Say what the reader should see.}
  \label{fig:myplot}
\end{figure}
\end{verbatim}

Draw figures in Python or with TikZ. Photographs of handwritten work are not
accepted.

\subsection*{Python code}
\begin{lstlisting}
def binary_search(a, t):
    lo, hi = 0, len(a) - 1
    while lo <= hi:
        mid = (lo + hi) // 2
        if a[mid] == t:
            return mid
        if a[mid] < t:
            lo = mid + 1
        else:
            hi = mid - 1
    return None
\end{lstlisting}

\subsection*{Pseudocode}
Tables, figures and pseudocode are \emph{floats}. LaTeX moves them to wherever
they fit best on the page, which may not be where you typed them. The
\texttt{[h]} asks for "here if possible" and is only a request. Do not fight
it, and never write "the table below". Refer to floats by number, as in
Table~\ref{tab:search}.

\begin{algorithm}[h]
\caption{Binary search}
\begin{algorithmic}[1]
  \State $lo \gets 1$, $hi \gets n$
  \While{$lo \le hi$}
    \State $mid \gets \lfloor (lo + hi)/2 \rfloor$
    \If{$A[mid] = t$} \State \Return $mid$ \EndIf
    \If{$A[mid] < t$} \State $lo \gets mid + 1$
    \Else \State $hi \gets mid - 1$ \EndIf
  \EndWhile
  \State \Return \textsc{Not Found}
\end{algorithmic}
\end{algorithm}

% =============================================================================
% END OF THE PART YOU DELETE
% =============================================================================


% --- The bibliography --------------------------------------------------------
% These two lines print the reference list. "apalike" is the citation style.
% "scribe-refs" is the bibliography file, without the .bib on the end. Only the
% sources you actually cited will appear.
\bibliographystyle{apalike}
\bibliography{scribe-refs}

\end{document}
